% Generated by tools/build.js from content/resume.yaml -- edit the YAML, not this file.
\documentclass[10pt,letterpaper]{article}
\usepackage[T1]{fontenc}
\usepackage[utf8]{inputenc}
\usepackage{lmodern}
\usepackage{textcomp}
\usepackage[margin=0.6in]{geometry}
\usepackage{enumitem}
\usepackage[hidelinks]{hyperref}
\hypersetup{pdftitle={Jesús M. Rueda-Becerril, Ph.D.}, pdfauthor={Jesús M. Rueda-Becerril}}
% map every glyph (including ligatures) to Unicode so text extraction is exact
\input{glyphtounicode}
\pdfgentounicode=1

\pagestyle{empty}
\setlength{\parindent}{0pt}
\setlength{\parskip}{0pt}
\setlength{\emergencystretch}{3em}
\setlist[itemize]{leftmargin=1.3em, topsep=1pt, itemsep=0.5pt, parsep=0pt}
\newcommand{\cvsection}[1]{%
  \vspace{5pt}{\large\bfseries #1}\par\nopagebreak\vspace{2pt}\hrule
  \nopagebreak\vspace{3pt}\nopagebreak}
\newcommand{\cvsubsection}[1]{%
  \vspace{3pt}{\bfseries #1}\par\nopagebreak\vspace{2pt}\nopagebreak}
\newcommand{\entrygap}{\vspace{3pt}}
\newcommand{\leftright}[2]{{#1\unskip\nobreak\hfil\penalty50\hskip1em\hbox{}%
  \nobreak\hfil\mbox{#2}\parfillskip=0pt\par}}

\begin{document}

\begin{center}
  {\LARGE\bfseries Jesús M. Rueda-Becerril, Ph.D.}\\[3pt]
  Seattle, WA \textbar{} \href{mailto:jm.ruebe@gmail.com}{jm.ruebe@gmail.com} \textbar{} +1 (765) 430-2330 \textbar{} \href{https://linkedin.com/in/jeruebe}{linkedin.com/in/jeruebe} \textbar{} \href{https://altjerue.github.io}{altjerue.github.io}
\end{center}

\cvsection{Professional Summary}
Computational physicist (Ph.D.) with 10+ years writing scientific code, running computational experiments, and validating models against data. Currently evaluates frontier language models on expert-level physics reasoning and builds neural-network surrogates of expensive simulations. Author of open-source scientific libraries (Paramo, Tleco; ApJ 2024) with code review and unit testing, and mentor of graduate students.\par

\cvsection{Skills}
\textbf{Languages:} Python, C/C++, Fortran 95, R, Rust, Octave, SQL, Java, Julia, Shell\par
\textbf{Simulation \& Modeling:} High-fidelity physics simulation (multiscale/multiphysics, kinetic plasma modeling, radiative processes (synchrotron, inverse Compton), radiative transfer, relativistic MHD, hydrodynamics (SPH), stochastic particle dynamics, Lagrangian particle dispersion), numerical PDEs (Fokker–Planck, advection–diffusion), finite difference/volume methods, numerical linear algebra, model validation, physics-informed synthetic data generation, uncertainty quantification, parameter estimation\par
\textbf{ML \& Data:} Scientific ML, language-model evaluation (scientific reasoning), surrogate modeling, Bayesian inference, Gaussian process regression, gradient-based optimization, PyTorch, TensorFlow, scikit-learn, NumPy, SciPy, Pandas, Astropy, neural networks, SVM, random forest, JAX (introductory coursework)\par
\textbf{HPC \& GPU:} OpenMP, MPI, OpenACC (GPU exploration), SLURM, PBS, HDF5, large-scale parallel computing (Frontera)\par
\textbf{Geospatial:} QGIS, raster/vector processing, Sentinel-1/2 multispectral data, DEM terrain analysis, terrain and elevation derivatives, geospatial feature engineering for ML\par
\textbf{Tools:} Git and code review, Docker, Azure, CI/CD (Jenkins, OpenShift), Kafka, Splunk, Jira, QGIS, ParaView, VisIt\par

\cvsection{Experience}
\leftright{\textbf{Independent Contractor}}{Jan 2026 – Present}\nopagebreak
\leftright{\textit{Mercor}}{Remote}\nopagebreak
Physics subject-matter expert evaluating frontier language models for multiple leading AI labs.\par\nopagebreak
\begin{itemize}
  \item{} Review 50+ expert-level physics, hydrodynamics, astrophysics, and cosmology problems for leading AI labs; write original prompts that stump frontier models and co-develop others with fellow writers.
  \item{} Diagnose failure modes in model reasoning and write structured feedback to improve scientific accuracy.
  \item{} Benchmark model performance and quality-check solutions to complex STEM problems.
\end{itemize}
\entrygap

\leftright{\textbf{Spatial Data Scientist}}{May 2025 – Jan 2026}\nopagebreak
\leftright{\textit{TealWaters}}{Seattle, WA}\nopagebreak
Built geospatial ML pipelines and optimized terrain-analysis code for wetland mapping.\par\nopagebreak
\begin{itemize}
  \item{} Mapped wetland probability across the 2,160 km² Skykomish watershed with the Wetland Intrinsic Potential (WIP) random-forest model, reaching >75\% accuracy.
  \item{} Optimized the terrain-modeling code that computes elevation derivatives from 1 m, 3 m, and 10 m DEMs over 20 tiles, improving its computational performance.
  \item{} Engineered features from Sentinel-1/2 multispectral imagery and terrain derivatives for ML land-cover and wetland models; built Python tools and QGIS visualizations.
  \item{} Presented results to managers and decision-makers, trained 2 colleagues on the WIP tool, and worked with software engineers to move prototype tools into production.
\end{itemize}
\entrygap

\leftright{\textbf{Independent Researcher \& Open-Source Developer}}{Jan 2024 – May 2025}\nopagebreak
\leftright{\textit{Self-directed}}{Seattle, WA}\nopagebreak
Continued simulation research and open-source development between industry roles.\par\nopagebreak
\begin{itemize}
  \item{} Co-developed Tleco, an open-source Rust/Python toolkit that solves the Fokker–Planck equation coupled to radiative transfer to model radiation from relativistic outflows (ApJ 2024).
  \item{} Developing WindsOfChange with the UW Forest Mycobiome Lab: a C++/R library that simulates wind-driven spore dispersal over forested, hilly terrain, parallelized with RcppParallel and OpenMP.
  \item{} Used AI coding assistants in the most recent stages of development; changes were accepted only if they passed the unit-test suite and, for some functions, matched reference solutions.
  \item{} Invited talk on machine learning in environmental science at the University of Washington (2024).
\end{itemize}
\entrygap

\leftright{\textbf{Software Engineer}}{Apr 2022 – Jan 2024}\nopagebreak
\leftright{\textit{Paychex}}{Remote}\nopagebreak
Built and tested data-integration services between enterprise systems and Oracle E-Business Suite.\par\nopagebreak
\begin{itemize}
  \item{} Built Java Kafka consumers that moved data between enterprise databases, deployed through Jenkins and OpenShift CI/CD pipelines.
  \item{} Created Splunk dashboards tracking the 200–400 users shared between enterprise systems and Oracle EBS, with alerts that flagged production incidents.
  \item{} Prepared, validated, and analyzed Oracle EBS data with SQL.
  \item{} Wrote full test suites for 5+ Python services in the CI pipeline, keeping production defect rates low.
\end{itemize}
\entrygap

\leftright{\textbf{Postdoctoral Research Associate}}{Feb 2021 – Apr 2022}\nopagebreak
\leftright{\textit{Rochester Institute of Technology}}{Rochester, NY}\nopagebreak
HPC simulations of supermassive black hole binaries in an NSF-funded project.\par\nopagebreak
\begin{itemize}
  \item{} Led a 7-person team upgrading PatchworkMHD, a C code for magnetohydrodynamic simulations of supermassive black hole binaries.
  \item{} Added black hole spin to the simulations with under 5\% runtime overhead.
  \item{} Ran and benchmarked large production simulations on Frontera (TACC), about 10,000 node-hours, and identified performance bottlenecks in the parallel code.
  \item{} Calibrated Paramo parameters against Fermi-LAT gamma-ray spectra with gradient-based optimization, comparing simulation output with data and updating the model automatically; supervised a graduate student on the same problem.
  \item{} Co-authored 2 papers from a multi-institution study of binary neutron star mergers and a third on blazar heating and cooling.
\end{itemize}
\entrygap

\leftright{\textbf{Postdoctoral Research Fellow}}{Oct 2018 – Nov 2020}\nopagebreak
\leftright{\textit{Purdue University}}{West Lafayette, IN}\nopagebreak
Radiative-transfer modeling of relativistic jets (blazars and gamma-ray burst afterglows).\par\nopagebreak
\begin{itemize}
  \item{} Designed and wrote Paramo, an open-source Fortran 95 library that solves the Fokker–Planck equation for relativistic particle distributions coupled to radiative transfer; optimized it with OpenMP and tested GPU offload with OpenACC.
  \item{} Paramo has been adopted by other research groups, including Combi \& Siegel (2023).
  \item{} Served as co-I on two NASA Fermi Guest Investigator grant proposals; the funded one (\$68K) supported the finding that blazar jets launch with similar energy per baryon, independently of their power.
  \item{} Added numerical-integration and OpenMP features to Paramo and built Python tools to compute energy losses and spectra of high-energy particles in gamma-ray burst afterglows.
  \item{} Mentored 3 graduate students (1 M.S., 2 Ph.D.), co-authoring papers with each.
  \item{} Wrote Python scripts for statistical modeling of the COVID-19 outbreak in Mexico with a multidisciplinary group, and co-wrote Spanish-language infographics and blog posts against misinformation.
\end{itemize}
\entrygap

\leftright{\textbf{Postdoctoral Research Fellow}}{Jan 2018 – Sep 2018}\nopagebreak
\leftright{\textit{Universidad Michoacana de San Nicolás de Hidalgo}}{Morelia, Mexico}\nopagebreak
\begin{itemize}
  \item{} Built a Python pipeline that turned GRTrans black hole simulation images into training data for an SVM classifier of black hole spin (grayscale conversion, intensity normalization, spatial feature extraction).
  \item{} Released an open-source Python tool for radiative-transfer spectra and light curves.
  \item{} Organized an HDF5 workshop to train graduate students in high-volume data storage.
\end{itemize}
\entrygap

\leftright{\textbf{Graduate Research Assistant}}{Oct 2011 – Jul 2017}\nopagebreak
\leftright{\textit{Universitat de València}}{Valencia, Spain}\nopagebreak
\begin{itemize}
  \item{} Automated runs of the C-SPEV radiative-transfer code with Shell and Python pipelines: 10–60 simulations per project and tens of GB of HDF5 output, with data-quality checks for downstream analysis.
  \item{} Studied how magnetization of shocked plasma shapes the radiation spectra of relativistic internal shocks, using the C-SPEV simulation code.
  \item{} Analyzed multi-wavelength observations from the NASA Fermi-LAT telescope and the Very Long Baseline Array (VLBA), including cross-matching sources across catalogs by sky coordinates and fitting log-parabola spectra.
  \item{} Fit C-SPEV simulation models to telescope data, identifying spectra and light-curve patterns that quantify plasma magnetization in blazars.
  \item{} Extended C-SPEV to compute discrete and continuous spectra from particle distributions of arbitrary shape, without affecting simulation runtime.
  \item{} Published 2 first-author papers in MNRAS.
\end{itemize}
\entrygap


\cvsection{Education}
\leftright{\textbf{Ph.D. in Physics}}{Oct 2011 – Jul 2017}\nopagebreak
\leftright{\textit{Universitat de València}}{Valencia, Spain}\nopagebreak
Excellent \textit{cum laude}\par\nopagebreak
Advisors: Prof. Miguel A. Aloy and Dr. Petar Mimica\par\nopagebreak
Thesis: \href{http://roderic.uv.es/handle/10550/60003}{Numerical treatment of radiation processes in the internal shocks of magnetized relativistic outflows}\par
\entrygap

\leftright{\textbf{M.Sc. in Physics}}{Aug 2009 – Sep 2011}\nopagebreak
\leftright{\textit{Universidad Michoacana de San Nicolás de Hidalgo}}{Morelia, Mexico}\nopagebreak
Advisor: Prof. José A. Cervera\par\nopagebreak
Thesis: \textit{Study of TOV stars with the SPH method} (wrote an SPH hydrodynamics code to study the stability of relativistic stars)\par
\entrygap

\leftright{\textbf{B.Sc. in Physics}}{Aug 2004 – Dec 2008}\nopagebreak
\leftright{\textit{Universidad Autónoma del Estado de México}}{Toluca, Mexico}\nopagebreak
Advisor: Prof. Francisco S. Guzmán\par\nopagebreak
Thesis: \textit{Numerical solution of null geodesics for the generation of gravitational lenses produced by spherically-symmetric and static spacetimes}\par\nopagebreak
Lic. Juan Josafat Pichardo Cruz Award\par
\entrygap


\cvsection{Software and Projects}
\leftright{\textbf{WindsOfChange}}{Aug 2024 – Present}\nopagebreak
\leftright{\textit{Lagrangian particle-dispersion library (C++/R)}}{\href{https://github.com/UW-ForestMycobiomeLab/WindSporesHills}{github.com/UW-ForestMycobiomeLab/WindSporesHills}}\nopagebreak
\begin{itemize}
  \item{} Translated an existing MATLAB code to R and C++, producing a library for Lagrangian particle dispersion over hilly terrain with forest canopy, driven by prescribed wind fields, based on analytical models (Massman \& Weil 1999; Finnigan \& Belcher 2004). Developed with UW collaborators; manuscript in preparation.
  \item{} Parallelized particle tracking with RcppParallel and OpenMP; designed for scalability to tens of thousands of particles across kilometer-scale domains.
  \item{} Used AI coding assistants in the most recent stages of development; changes were accepted only if they passed the unit-test suite and, for some functions, matched reference solutions.
\end{itemize}
\entrygap

\leftright{\textbf{Blazar Population ML}}{Jan 2026 – Present}\nopagebreak
\leftright{\textit{Physics-informed ML (Python, scikit-learn, PyTorch), developed with Claude}}{\href{https://github.com/altjerue/ilwikak-ml}{github.com/altjerue/ilwikak-ml}}\nopagebreak
\begin{itemize}
  \item{} Applying classical ML (PCA, UMAP, Gaussian mixture models, HDBSCAN) to the Fermi-LAT 4LAC-DR3 catalog (3,407 sources) to test whether physical jet properties can be recovered from multi-wavelength observables; documented missingness and selection effects feature by feature.
  \item{} Fitting Gaussian process regression to a lower bound on the jet Lorentz factor for a 334-source MOJAVE cross-matched subsample.
  \item{} Building a neural-network surrogate of the forward model (jet parameters to observable SEDs), to be benchmarked against Tleco simulations, to enable fast parameter inference.
  \item{} Forecasting blazar gamma-ray flares with a time-series foundation model fine-tuned on Tleco-generated synthetic light curves (in progress).
\end{itemize}
\entrygap

\leftright{\textbf{Tleco}}{Jan 2024 – Sep 2024}\nopagebreak
\leftright{\textit{Radiative-transfer toolkit (Rust/Python)}}{\href{https://github.com/zkdavis/Tleco}{github.com/zkdavis/Tleco}}\nopagebreak
\begin{itemize}
  \item{} Numerical PDE solvers, particle-distribution algorithms, and spectral calculations for relativistic plasmas; successor to Paramo.
\end{itemize}
\entrygap

\leftright{\textbf{Paramo}}{Oct 2018 – Apr 2022}\nopagebreak
\leftright{\textit{Radiative-transfer library (Fortran 95, OpenMP)}}{\href{https://github.com/altjerue/paramo}{github.com/altjerue/paramo}}\nopagebreak
\begin{itemize}
  \item{} Fokker–Planck and radiative-transfer solver for relativistic particle distributions, parallelized with OpenMP; used by other groups, including Combi \& Siegel (2023).
\end{itemize}
\entrygap

\leftright{\textbf{PatchworkMHD}}{Feb 2021 – Apr 2022}\nopagebreak
\textit{Large-scale MHD simulation code (NSF project)}\par\nopagebreak
\begin{itemize}
  \item{} Led a 7-person upgrade of this production C code for HPC environments, adding black hole spin with under 5\% runtime overhead.
\end{itemize}
\entrygap


\cvsection{Selected Publications (9 published articles, 3 first-author; 3 in preparation)}
\begin{itemize}
  \item{} Davis, Z., \textbf{Rueda-Becerril, J. M.}, \& Giannios, D. \textit{Tleco: A Toolkit for Modeling Radiative Signatures from Relativistic Outflows}, \href{https://doi.org/10.3847/1538-4357/ad8bc2}{ApJ 976, 182 (2024)}, \href{https://arxiv.org/abs/2405.17581}{arXiv:2405.17581}.
  \item{} Davis, Z., \textbf{Rueda-Becerril, J. M.}, \& Giannios, D. \textit{Balancing Turbulent Heating with Radiative Cooling in Blazars}, \href{https://doi.org/10.1093/mnras/stac1282}{MNRAS 513, 5766–5779 (2022)}, \href{https://arxiv.org/abs/2201.07790}{arXiv:2201.07790}.
  \item{} \textbf{Rueda-Becerril, J. M.}, Harrison, A. O., \& Giannios, D. \textit{Blazar jets launched with similar energy per baryon, independently of their power}, \href{https://doi.org/10.1093/mnras/staa3925}{MNRAS 501, 4092–4102 (2021)}, \href{https://arxiv.org/abs/2009.02273}{arXiv:2009.02273}.
  \item{} \textbf{Rueda-Becerril, J. M.}, Mimica, P., \& Aloy, M. A. \textit{On the influence of a hybrid thermal–non-thermal distribution in the internal shocks model for blazars}, \href{https://doi.org/10.1093/mnras/stx476}{MNRAS 468, 1169–1182 (2017)}, \href{https://arxiv.org/abs/1612.06383}{arXiv:1612.06383}.
  \item{} Guzmán, F. S., \& \textbf{Rueda-Becerril, J. M.} \textit{Spherical boson stars as black hole mimickers}, \href{https://doi.org/10.1103/PhysRevD.80.084023}{Phys. Rev. D 80, 084023 (2009)}, \href{https://arxiv.org/abs/1009.1250}{arXiv:1009.1250}.
\end{itemize}

\end{document}
